% Faithful transcription of official CCC2024 article4, pp.4:2,4:8--4:9,4:15--4:17.
\paragraph{Theorem 3.} For every $\alpha\in[1,1.5]$,
\[\mathsf{BPL}\subseteq\mathsf{CTISP}[2^{\log^\alpha(n)},\log^\alpha(n),\log^{3-\alpha}(n)].\]
\section*{Author Section 2.1: Composition of Catalytic Algorithms}
We state the main result of this section, which is that catalytic algorithms can be composed without increasing the catalytic space usage. We must be careful when dealing with partial catalytic machines, and in this case we only obtain composition if the machines are reversible (Definition 11).
\paragraph{Theorem 12 (Composition of Partial Catalytic Machines).} Suppose reversible catalytic machines $M_1,M_2$ partially compute $f_1,f_2:\{0,1\}^n\to\{0,1\}^n$ respectively using workspace $S(n)\geq\log(n)$, catalytic space $C(n)$, and time $T(n)$. Then there is a reversible catalytic machine $M$ that partially computes $f_2\circ f_1$ using workspace $2S(n)+O(\log(Sn))$, catalytic space $C$, and time $\operatorname{poly}(T(n))$. Moreover, $M^{\mathbf w}(x)=\bot$ only if $M_2^{\mathbf w}(f_1(x))=\bot$, or there exists $\mathbf w'$ such that $M_1^{\mathbf w'}(x)=\bot$.
\paragraph{Proof.} We proceed roughly following the standard proof for composition of space-bounded algorithms. We maintain two sections on the worktape of size $S$ for $M_1$ and $M_2$ (and auxiliary state of size $O(\log(Sn))$ to keep track of the location of read and write heads, and the FSM configuration of both machines), and a single catalytic tape $\mathbf w$.
\paragraph{The Simulation.} We now begin to simulate $M_2^{\mathbf w}$. We first verify that $M_1^{\mathbf w}(x)\neq\bot$. As in the conventional composition of space-bounded algorithms, every time $M_2$ reads its input, we run $M_1^{\mathbf w}(x)$ on a separate section of the worktape and return the relevant bit of its output, where $\mathbf w$ is the same catalytic tape used by $M_2$, in whatever its current configuration is at the time of the tape read. Moreover, every time $M_2$ writes to the catalytic tape, resulting in a configuration $\mathbf w'$, we run $M_1^{\mathbf w'}(x)$ and verify that it does not produce $\bot$.
\paragraph{Computing the Function.} In the case that $M_1^{\mathbf w'}(x)=f_1(x)$ for every configuration $\mathbf w'$ that is encountered in this simulation and $M_2^{\mathbf w}(f_1(x))=f_2(f_1(x))$, it is easy to see that $M^{\mathbf w}(x)=f_2(f_1(x))$. Moreover, it is clear that in this case we successfully reset the tape.
Otherwise, consider the first point at which $M_1^{\mathbf w'}(x)=\bot$. We first undo the most recent change to the catalytic tape, and send the REVERT command to $M_2$. Once $M_2$ has finished reverting, return $\bot$. We claim that $M_2$ successfully reverts the tape. This follows from the reversibility of $M_2$, and the fact that all calls $M_2$ makes to its input during this process are correctly answered by $M_1$. The latter property follows as every time $M_2$ queries its input during the revert process, $\mathbf w$ is in a state that was encountered during the forward pass, and hence $M_1(x)$ produced $f_1(x)$ when initialized with this catalytic configuration (as otherwise we would have aborted sooner).
\paragraph{Reversibility.} Essentially the same argument establishes that $M$ is reversible. If we receive the REVERT command, let $P(\mathbf w)$ be the states of the catalytic tape that have been encountered so far. First send REVERT to $M_1$ (if operating), and once it has completed send the REVERT command to $M_2$. Moreover, while $M_2$ is reverting, we claim that $\mathbf w$ remains in $P(\mathbf w)$. This follows from the fact that $M_2$ is reversible (as any configurations it creates will lie in $P(\mathbf w)$), and moreover every time $M_2$ queries its input, any computation done by $M_1$ will likewise keep $\mathbf w$ in $P(\mathbf w)$, as we already called $M_1$ with this starting configuration in the forward pass (and $M_1$ will not produce $\bot$, as otherwise we would have already aborted). Thus, the composed algorithm is reversible.
\paragraph{Time and Space.} We now argue the space and time are as claimed. There are a constant number of pointers (which we maintain on the worktape) to track the number of bits output by $M_1$, current tape heads, and other information. The fact that the catalytic tape size is preserved is immediate. The call overhead adds at most a polynomial factor in the runtime, as we run $M_1$ at most once per step of $M_2$. Finally, if $M_1$ computes $f_1(x)$ for every catalytic tape and $M_2$ computes correctly on $\mathbf w$, we successfully compute $f_2\circ f_1$ as claimed.\hfill$\triangleleft$
We derive an easy corollary in the case of multiple composition:
\paragraph{Corollary 13.} Suppose a reversible catalytic machine $M$ partially computes $f:\{0,1\}^n\to\{0,1\}^n$ using workspace $S(n)\geq\log n$, catalytic space $C$, and time $2^S$. Then there exists a reversible catalytic machine $M'$ that partially computes $f^\ell$ using workspace $O(\ell\cdot S)$, catalytic space $C$, and time $2^{O(\ell\cdot S)}$. Moreover, for $x$ where $M^{\mathbf w}(f^i(x))=f(f^i(x))$ for every $\mathbf w$ and $i\in\{0,\ldots,\ell-1\}$, $M'^{\mathbf w}(x)=f^\ell(x)$.
\section*{Author Section 4: Catalytic Recursive Matrix Powering}
We now transform Theorem 23 into a parameterized algorithm for matrix powering.
\paragraph{Theorem 27.} There is a catalytic machine that, given $r_1,r_2$ such that $r_1r_2=\log(n)$ and a stochastic matrix $M\in[0,1]^{n\times n}$ where each entry is specified with $l=O(\log n)$ bits of precision, uses workspace $O(r_2\cdot\log(n))$, catalytic space $O(r_1\cdot\log(n))$, and time $n^{O(r_2)}$, and outputs $\widetilde M$ such that $\|\widetilde M-M^n\|\leq1/n$.
Theorem 27 immediately implies Theorem 3 by setting $r_2=\log^{\alpha-1}(n)$ and $r_1=\log^{2-\alpha}(n)$ for $\alpha\in[1,1.5]$, and using the standard transformation of estimating the acceptance probability of a BPL machine via stochastic matrix powering.
We first prove there exists an algorithm which computes a single intermediate power. We must be careful to ensure that the algorithm satisfies the requirements of (partial) catalytic computation. In particular, if the machine ever outputs an answer (rather than $\bot$), this must be the only possible answer for this input, over all possible catalytic tapes. Simultaneously, we must ensure that the vast majority of inputs never return $\bot$ no matter the initial catalytic tape configuration.
We achieve this dual guarantee using an idea from Saks and Zhou [21]. For a given input matrix $M$, we additionally take in a shift $s\in[0,\delta]$ for $\delta=1/\operatorname{poly}(w/\varepsilon)$. After computing an approximate power of $M$, we add $s$ to each entry, and then truncate each entry to $O(\log(w/\varepsilon))$ bits of precision. In fact, we first verify that our shifted approximate power is sufficiently far from the rounding threshold, and if not return $\bot$. By doing so, we algorithmically verify that we will never round to different thresholds over different $\mathbf w$. Unfortunately, for some pairs $(M,s)$ it may be the case that we detect possible mis-rounding for some tapes $\mathbf w$, even if all possible approximations lie inside a single rounding interval. This can result in returning $\bot$ on some tapes and a (consistent) value otherwise. However, we show that we can choose the magnitude of $s$ such that with high probability over $s$ this does not occur, and we always return a (consistent) value.
\paragraph{Theorem 28.} For every $n,w\in\mathbb N$ and $\varepsilon>0$ where $w\geq n\geq\log w$ and $2^{-n}>\varepsilon>0$, there is a reversible catalytic machine $P$ that uses workspace $O(\log(w/\varepsilon))$, catalytic space $O(\log(n)\log(nw))$, and time $\operatorname{poly}(nw/\varepsilon)$. The machine takes input $s\in\{0,1\}^{O(\log(w/\varepsilon))}$ and a substochastic matrix $M\in[0,1]^{w\times w}$, where each entry of $M$ is specified with $l=O(\log(w/\varepsilon))$ bits of precision. Moreover:
\begin{itemize}
\item For every $(s,M)$, there is substochastic $\widetilde M_s$ (defined without reference to $\mathbf w$) with $l$ bits of precision satisfying $\|\widetilde M_s-M^n\|_1\leq\varepsilon$ such that for every $\mathbf w$, $P^{\mathbf w}(s,M)\in\{\bot,\widetilde M_s\}$.
\item For every $M$, $\Pr_s[\exists\mathbf w,P^{\mathbf w}(s,M)=\bot]\leq1/w^2$.
\end{itemize}
\paragraph{Proof.} Let $B:[(w/\delta)^c]\times\{0,1\}^{m=n\cdot O(\log(w/\delta))}\to[(w/\delta)^c]$ be the result of Lemma 15 applied to $M$ with $n=n$ and $\varepsilon=\delta$ to be chosen later. We compose the output of this algorithm with Theorem 23, applied with $B=B$ and $w'=(w/\delta)^c$ and $n'=m=O(n^3)$ and $\varepsilon=\delta$ to be chosen later. Let $\vec h=(h_1,\ldots,h_\ell)$ be the hash functions obtained from this call, which we have implicit access to via the current state of the catalytic tape $\mathbf w'$ and $V$, and let $t=O(\log(w/\delta))$ be the domain of the hash functions. Then enumerate over $x\in\{0,1\}^t$ and for every $i,j\in[w]$ let
\[\widetilde M_{i,j}=\Pr_{x\leftarrow U_t}[B[i,\operatorname{NIS}_{\vec h}(x)]=j].\]
Next, define $\tau=s\cdot2^{-2k}\cdot J$ where $J=1^{w\times w}$, interpreting $s\in[2^k]$. We next check if any entry of $(\widetilde M+\tau)$ is within $2\delta$ of a multiple of $2^{-l}$, our rounding threshold. In this case, run $D^{\mathbf w'}(B,V)$ to reset the tape and return $\bot$. Otherwise, let
\[\widetilde M_s=\lfloor\widetilde M+\tau\rfloor_l\]
where $\lfloor\cdot\rfloor_l$ rounds each entry down to $l$ bits of precision, and decreases the largest entry per row such that the final matrix is substochastic. Let this matrix be $\widetilde M_s$. Finally, run $D^{\mathbf w'}(B,V)$, and return $\widetilde M_s$.
\paragraph{Accuracy.} By our choice of error in Theorem 23 and Lemma 15, we have that
\[\|\widetilde M-M^n\|_1\leq2w\delta\]
and moreover $\widetilde M$ has each row sum at most 1. Furthermore, perturbing by $\tau$ and rounding down the largest entry causes an $\ell_1$ error of at most $2w\cdot2^{-k}$. Finally, rounding each entry down to a multiple of $2^{-l}$ causes a total error of at most $w\cdot2^{-l}$, so
\[\|\widetilde M_s-M^n\|\leq2w\delta+2w\cdot2^{-k}+w\cdot2^{-l}\leq\varepsilon\]
Where the final inequality comes from choosing
\[l=O(\log(w/\varepsilon)),\qquad k=10\cdot l,\qquad\delta=2^{-2k}.\]
\paragraph{Uniqueness.} We claim that for every $(s,M)$, there is at most 1 possible non-$\bot$ output over all choices of $\mathbf w$ (which we denote $\widetilde M_s$ in the theorem statement). Let $\widetilde M^{\mathbf w},\widetilde M^{\mathbf w'}$ be the result of Theorem 23 on $M$ initialized with catalytic tapes $\mathbf w,\mathbf w'$. By the accuracy guarantee of Theorem 23, for every $i,j$ we have
\begin{align*}
|(\widetilde M^{\mathbf w}+\tau)_{i,j}-(\widetilde M^{\mathbf w'}+\tau)_{i,j}|&\leq|(\widetilde M^{\mathbf w}+\tau)_{i,j}-(M+\tau)_{i,j}|\\
&\quad+|(M+\tau)_{i,j}-(\widetilde M^{\mathbf w'}+\tau)_{i,j}|\leq2\delta.
\end{align*}
Thus, if $(\widetilde M^{\mathbf w'}+\tau)_{i,j}$ is greater than $2\delta$ from a multiple of $2^{-l}$, we can be certain that no tape $\mathbf w'$ will induce an estimate that falls on the other size of the threshold, and hence all non-$\bot$ outputs will be rounded consistently.
\paragraph{Success Probability.} Furthermore, we argue that for every $M$, with probability at least $1-1/w^2$ over $s$ we return $\widetilde M_s$ (not $\bot$) for every initial tape configuration. Fixing arbitrary $s$ and $i,j\in[w]$, if $(M+\tau)_{i,j}$ is at least $3\delta$ from every multiple of $2^{-l}$, every $\mathbf w$ will induce an estimate $(\widetilde M+\tau)_{i,j}$ that is at least $2\delta$ from every multiple of $2^{-l}$, and hence for every $\mathbf w$ we will not produce $\bot$ due to this entry. This occurs for every $i,j$ simultaneously with probability at least
\[w^2\cdot2^l\cdot\frac{6\delta\cdot2^k+2}{2^k}\ll1/w^2.\]
\paragraph{Time and Space.} It is clear the algorithm runs in the claimed time and space bound, given Theorem 23.
\paragraph{Reversibility.} As the only components of the algorithm that write to the catalytic tape are calls to Theorem 23, reversibility follows immediately from the equivalent result for that algorithm.\hfill$\triangleleft$
We can then prove the main result.
\paragraph{Theorem 27 (restated).} There is a catalytic machine that, given $r_1,r_2$ such that $r_1r_2=\log(n)$ and a stochastic matrix $M\in[0,1]^{n\times n}$ where each entry is specified with $l=O(\log n)$ bits of precision, uses workspace $O(r_2\cdot\log(n))$, catalytic space $O(r_1\cdot\log(n))$, and time $n^{O(r_2)}$, and outputs $\widetilde M$ such that $\|\widetilde M-M^n\|\leq1/n$.
\paragraph{Proof.} Let $\vec s=(s_1,\ldots,s_{r_2})\in\{0,1\}^{r_2\cdot O(\log n)}$ be a vector of random shifts. Let $\widetilde M_0=M$ and for $i\in[r_2]$ recursively define
\[\widetilde M_i=f(\widetilde M_{i-1},s_i),\]
where $f$ is the function defined by Theorem 28 with $\varepsilon=1/n^3$ and $n=2^{r_1}$. An easy inductive proof [19] establishes that, letting
\[\|\widetilde M_i-M^{2^{r_1\cdot i}}\|_1=\delta_i\]
we have $\delta_{i+1}\leq1/n^3+2^{r_1}\cdot\delta_{i-1}\leq2^{r_1+1}\cdot\delta_{i-1}$, and hence $\delta_{r_2}\leq1/n$.
The final algorithm iterates over $\vec s$ and computes $\widetilde M_{r_2}$ by applying recursive composition of space-bounded machines Corollary 13 to the algorithm of Theorem 28 as defined above.\footnote{We can define the machine to take the entire shift vector and a pointer to the index it should use, such that we are recursively composing exactly the same machine, but we suppress this for clarity.} The algorithm returns the first non-$\bot$ output. The fact that the algorithm is catalytic follows from Corollary 13 and Theorem 28. Next, we claim there is some $\vec s$ where the algorithm returns a value. Note that $s_i$ is chosen obliviously to $\widetilde M_{i-1}$, and so with probability at least $1/n^2$ over $s_i$, on input $(\widetilde M_{i-1},s_i)$ the algorithm returns $\widetilde M_i$ (i.e. not $\bot$) when run with every possible catalytic tape. Thus, there is some $\vec s$ where every level computes correctly.
\paragraph{Time and Space.} Every application of Theorem 28 occurs with parameters $n=2^{r_1}$ and $w=n$ and $\varepsilon=1/n^3$, such that the algorithm uses workspace $O(r_1+\log(n))=O(\log n)$, catalytic space $O(r_1\cdot(r_1+\log(n)))=O(r_1\cdot\log n)$, and time $\operatorname{poly}(n)$, and moreover the shift $s$ for each level is of length $O(\log n)$. Applying Corollary 13, we obtain that the composed algorithm uses workspace $O(r_2\cdot\log(n)+|\vec s|)=O(r_2\cdot\log n)$, catalytic space $O(r_1\cdot\log n)$, and runs in time $n^{O(r_2)}$ as claimed.\hfill$\triangleleft$
